%% annex/p07/interpolation-supplement.tex
%% Complete Maehara, relational adequacy, general Craig, and Beth proofs.

\section{Maehara's partition method}

\begin{annexitem}{annex:p7.interp.1}{thm:p07-maehara-partition}{Maehara's partition lemma}
\begin{proof}
Induct on the height of the supplied cut-free
\(\mathsf{LK}_{\mathrm d}\)-derivation.  For an induced premise partition,
write \(I\), or \(I_0,I_1\), for the interpolants supplied by the induction
hypothesis.  We check every rule on the selected sheet.

For the initial sequent \(A\Rightarrow A\), there are four placements of its
two occurrences.  If both lie in the first half, take \(I=\bot\); the first
required sequent is identity with right weakening and the second is
\(\bot L\).  If both lie in the second half, take
\(I=\top\), where \(\top\) abbreviates \(\bot\to\bot\); derive
\(\Rightarrow\top\) by \(\to R\) over identity and use weakening with the
second-half identity.  If the antecedent occurrence is in the first half and
the succedent occurrence in the second, take \(I=A\) and use identity twice.
If the antecedent occurrence is in the second half and the succedent occurrence
in the first, take \(I=\neg A\), that is \(A\to\bot\): \(\to R\) over
contexted identity gives \(\Rightarrow A,\neg A\), and \(\to L\) from
\(A\Rightarrow A\) and \(\bot\Rightarrow\) gives
\(\neg A,A\Rightarrow\).  Contexts are restored by weakening.  In the two
same-half cases the interpolant is logical; in each cross case every symbol of
\(A\) occurs in both halves.  The free-variable condition is immediate.

For \(\bot,\Gamma\Rightarrow\Delta\), take \(I=\bot\) when the displayed
\(\bot\) lies in the first antecedent part and \(I=\top\) when it lies in the
second.  The required sequent containing \(\bot\) is initial, and the other is
obtained from \(\bot\Rightarrow\) or \(\Rightarrow\top\) by weakening.

Left or right weakening retains the premise interpolant and applies the same
weakening to whichever required sequent receives the new occurrence.
For contraction, lift the conclusion partition to the premise by putting both
contracted copies in the same half as the surviving occurrence; retain the
premise interpolant and contract in the corresponding required sequent.
Exchange requires no case because contexts are finite multisets.

The one-premise propositional rules are
\(\land L\), \(\lor R\), and \(\to R\).  Put every immediate subformula
displayed in the premise in the same half as the principal formula in the
conclusion.  If that is the first half, apply the original rule to the first
interpolant sequent and leave the second unchanged; if it is the second half,
leave the first unchanged and apply the rule to the second.  This produces
both required conclusion sequents with the same \(I\).  No formula outside the
principal formula is introduced, so the symbol and free-variable conditions
are preserved.

The two-premise propositional rules are
\(\land R\), \(\lor L\), and \(\to L\).  Their premise partitions retain the
half assignment of every context occurrence; for \(\to L\), this applies
separately to the two split context pairs.  Suppose first that the principal
formula is in the first half.  Put \(I=I_0\lor I_1\).  In each first-half
premise sequent, right-weaken and use \(\lor R\) to replace \(I_i\) by
\(I_0\lor I_1\), then apply the original two-premise rule.  On the second
half, align contexts by weakening and apply \(\lor L\) to the two sequents
beginning with \(I_0\) and \(I_1\).  If the principal formula is in the
second half, put \(I=I_0\land I_1\).  Align the first-half contexts and apply
\(\land R\).  For each second-half premise, left-weaken by the unused
\(I_j\) and apply \(\land L\), then apply the original two-premise rule.
When the split-context \(\to L\) rule concatenates two copies of the chosen
interpolant, apply the matching contraction; all other added context is aligned
by weakening before the rule.
These constructions cover both placements of all three branching rules.  A
symbol or free variable in \(I_0\lor I_1\) or \(I_0\land I_1\) was common in
one induced premise partition and remains present on both corresponding halves
of the conclusion; hence both side conditions persist.

In the quantifier cases, occurrence tests on a variable in an interpolant
refer to its free occurrences.  Choose each new bound variable \(z\)
fresh; \(I(z)\) denotes capture-avoiding replacement of the designated
free variable \(a\) by \(z\) in \(I(a)\).

For \(\forall L\) and \(\exists R\), the displayed term is a variable \(a\)
because the signature is purely relational.  Put \(A(a)\) in the same half as
the principal formula.  If \(a\) is absent from \(I\), or remains free in both
halves of the conclusion partition, retain \(I\) and apply the original rule
on the principal half.  The only remaining case is that \(a\) occurs in
\(I(a)\) but ceases to be common when \(A(a)\) is replaced by the principal
quantified formula.

If that principal formula is in the first half, put
\(J=\forall zI(z)\).  On the first interpolant sequent apply the original
\(\forall L\) or \(\exists R\) rule first, thereby removing the only
first-half occurrence of \(a\), and then apply \(\forall R\) to \(I(a)\).
On the second interpolant sequent apply \(\forall L\) to \(J\), instantiating
it with \(a\), and then use the induction-hypothesis derivation.  If the
principal formula is in the second half, put \(J=\exists zI(z)\).  The first
interpolant sequent obtains \(J\) by \(\exists R\).  In the second, uniformly
rename the instantiation variable to a fresh \(b\), apply the original rule to
remove \(A(b)\), and then apply \(\exists L\) to \(J\).  Its eigenvariable
condition holds because \(b\) is absent from the conclusion partition.
These are both placements of both non-eigenvariable quantifier rules.  They
remove every variable that the premise partition, but not the conclusion
partition, made common; no nonlogical function or constant symbol exists to
require a further lifting step.

Consider \(\forall R\), with eigenvariable \(a\) renamed fresh for the entire
conclusion partition.  If \(\forall xA\) lies in \(\Delta_1\), the induction
hypothesis has the form
\[
 \Gamma_1\Rightarrow\Delta'_1,A(a),I(a),
 \qquad
 I(a),\Gamma_2\Rightarrow\Delta_2.
\]
If \(a\) is absent from \(I\), retain \(I\) and apply \(\forall R\) to the
first sequent.  Otherwise put \(J=\exists zI(z)\), using a fresh bound
variable \(z\).  Apply \(\exists R\) to the occurrence of \(I(a)\) in the
first sequent and then \(\forall R\) to \(A(a)\); apply \(\exists L\) to the
second sequent.  Freshness holds because \(a\) is absent from the conclusion
and becomes bound in \(J\).  If \(\forall xA\) lies in \(\Delta_2\), the
induction-hypothesis sequents instead place \(A(a)\) in the second one.  When
needed put \(J=\forall zI(z)\): apply \(\forall R\) to the first interpolant
sequent, and in the second apply \(\forall L\) to recover \(I(a)\) before
applying \(\forall R\) to \(A(a)\).

The \(\exists L\) case is dual but we record it explicitly.  If
\(\exists xA\) lies in \(\Gamma_1\), the premise interpolant occurs in
\[
 A(a),\Gamma'_1\Rightarrow\Delta_1,I(a),
 \qquad
 I(a),\Gamma_2\Rightarrow\Delta_2.
\]
When \(a\) occurs in \(I\), put \(J=\exists zI(z)\); use \(\exists R\) on
\(I(a)\) and then \(\exists L\) on \(A(a)\) in the first sequent, and use
\(\exists L\) on the second.  If the principal formula lies in
\(\Gamma_2\), put \(J=\forall zI(z)\); use \(\forall R\) in the first
sequent, and in the second use \(\forall L\) to recover \(I(a)\) before
\(\exists L\) on \(A(a)\).  If \(a\) is absent from \(I\), retain \(I\)
and apply only the original eigenvariable rule on the appropriate half.

Quantifying \(a\) introduces no nonlogical symbol and removes the only free
variable not present in the conclusion partition.  All four eigenvariable
placements therefore satisfy both interpolation side conditions.  Cut and Mix
do not occur by hypothesis, and the selected sheet has no equality or
primitive-negation rules.  The cases above exhaust the calculus and complete
the induction.
\end{proof}
\end{annexitem}

\section{Relational cut-free completeness}

\begin{annexitem}{annex:p7.interp.2}{thm:p07-relational-cut-free-completeness}{Relational soundness and cut-free completeness}
\begin{proof}
Use exactly the rules of \texttt{sheet:lk-defneg}, with variables as the
only terms and \(\neg A\) expanded as \(A\to\bot\).

\paragraph{Soundness.}
Call a structure and assignment a counterexample to
\(\Gamma\Rightarrow\Delta\) when every member of \(\Gamma\) is true
and every member of \(\Delta\) is false.
Identity and \(\bot L\) have no counterexamples.
Weakening can only add a required true antecedent or false succedent,
and contraction removes repeated copies of one unchanged truth condition.
Consequently a counterexample to the conclusion of a structural rule
is a counterexample to its premise.

For \(\land L\), a true conjunction makes both premise conjuncts true;
for \(\lor R\), a false disjunction makes both disjuncts false.
A false implication in \(\to R\) has a true antecedent and false
consequent.  These observations give a counterexample to the unique
premise in each case.
For \(\land R\), a false principal conjunction has a false conjunct,
selecting a false premise.  For \(\lor L\), a true principal disjunction
has a true disjunct, again selecting a false premise.
For split-context \(\to L\), a true principal \(A\to B\) has either
false \(A\) or true \(B\).  The former falsifies the premise with
\(A\) in its succedent, and the latter falsifies the premise with
\(B\) in its antecedent.  The counterexample already satisfies both
split antecedent contexts and falsifies both split succedent contexts.

For \(\forall L\), truth of \(\forall xA\) gives truth of the displayed
instance \(A[a/x]\).  For \(\exists R\), falsity of \(\exists xA\)
gives falsity of every such instance.  Both use the substitution lemma.
If the conclusion of \(\forall R\) is false, choose a value falsifying
its quantified matrix and assign it to the fresh eigenvariable.
The other conclusion formulas keep their truth values by freshness,
so the premise is false.  If the conclusion of \(\exists L\) is false,
choose a witnessing value for its true existential antecedent and make
the same fresh-variable assignment; the premise is false.
Finally, a counterexample to a Cut conclusion falsifies the first
premise when the cut formula is false, and the second when it is true.
The same argument handles all displayed copies in Mix.
Induction up a derivation proves soundness.

\paragraph{Persistent backward steps.}
For completeness it suffices to consider the finite predicate signature
occurring in the endsequent.  A model of that signature can interpret
every unused predicate as empty.  Fix a countably infinite set
\(V=\{v_0,v_1,\ldots\}\) of variables, containing all variables of the
endsequent.
In this proof, \emph{unprovable} means having no derivation without Cut or Mix.
Suppose the endsequent is unprovable.
We construct larger finite antecedent and succedent sets, retaining all
formulas already present.  Using sets during this construction is harmless:
the registered weakening and contraction rules convert them to and from
the corresponding finite multiset contexts.

The following backward steps add the indicated formulas.  A vertical
choice denotes two alternative successors.
\[
\begin{array}{c|c|l}
 \text{side}&\text{principal formula}&\text{added formulas}\\ \hline
 \Gamma&A\land B&A,B\text{ in }\Gamma\\
 \Delta&A\lor B&A,B\text{ in }\Delta\\
 \Delta&A\to B&A\text{ in }\Gamma,\ B\text{ in }\Delta\\
 \Gamma&A\lor B&A\text{ in }\Gamma\ \mid\ B\text{ in }\Gamma\\
 \Delta&A\land B&A\text{ in }\Delta\ \mid\ B\text{ in }\Delta\\
 \Gamma&A\to B&A\text{ in }\Delta\ \mid\ B\text{ in }\Gamma\\
 \Gamma&\forall xA&A[a/x]\text{ in }\Gamma,\quad a\in V\\
 \Delta&\exists xA&A[a/x]\text{ in }\Delta,\quad a\in V\\
 \Gamma&\exists xA&A[b/x]\text{ in }\Gamma,\quad b\text{ fresh}\\
 \Delta&\forall xA&A[b/x]\text{ in }\Delta,\quad b\text{ fresh}
\end{array}
\]
All substitutions are capture avoiding.
For either eigenvariable step, choose \(b\) absent from the entire
current sequent and from the earlier eigenvariable choices.
Only finitely many variables are excluded at a finite stage.

If every successor of one of these steps were provable, the current
sequent would be provable.  To see this, apply the corresponding
registered logical rule to the successor derivation or derivations,
then contract the extra copy of the retained principal formula.
For \(\forall L\) and \(\exists R\), the chosen \(a\) is an allowed
term, so the same simulation applies to every \(a\in V\).
For the two eigenvariable rules, the displayed freshness condition
is exactly the one required by the forward inference.
For split \(\to L\), the two successor derivations have the whole
current context.  Applying \(\to L\) duplicates that context on both
sides; contract these copies and the repeated principal implication.
No Cut or Mix is used in any of these simulations.
An identity overlap, or \(\bot\) in the antecedent, would also give
an immediate proof by identity and weakening, or by \(\bot L\).
Thus an unprovable current sequent has no such closure and has at
least one unprovable successor at every applicable step.

\paragraph{Fair saturation.}
There are countably many formulas in the finite predicate signature
over \(V\).  Enumerate the possible expansion tasks, including a separate
task for every pair consisting of a \(\forall\)-antecedent or
\(\exists\)-succedent formula and a variable \(a\in V\).
In pass \(s\), inspect the first \(s\) tasks.  Perform a task when
its principal formula is present and the task has not already been
performed.  For a branching task, choose an unprovable successor;
for an eigenvariable task, record its fresh witness and perform it
once.  Every pass is finite.  An inapplicable task remains available
in later passes, and a pass with no applicable tasks leaves the
sequent unchanged.

This is a classical existence construction: it does not presume an
algorithm deciding whether a sequent is unprovable.  There are at
most two successor choices, and the preceding paragraph proves that
an unprovable one exists.  Recursion, selecting the first such choice
in a fixed order, gives an increasing sequence
\((\Gamma_s,\Delta_s)\) of finite unprovable sequents.
Let
\[
 \Gamma_\infty=\bigcup_s\Gamma_s,\qquad
 \Delta_\infty=\bigcup_s\Delta_s.
\]
Every task for a formula that enters either union is inspected in
every sufficiently late pass, so it is eventually performed.
In particular, every \(\forall\)-antecedent and
\(\exists\)-succedent receives its instance for \emph{every} variable
in \(V\), including variables not free in the initial sequent.
Each \(\exists\)-antecedent and \(\forall\)-succedent has a witnessing
instance in the required union.  All six propositional saturation
conditions in the table hold.
No formula belongs to both unions: its two occurrences would already
coexist in one finite stage and close that sequent.
Also \(\bot\notin\Gamma_\infty\).

\paragraph{The countermodel.}
Take domain \(V\), and let the assignment \(s_0\) send each variable
to itself.  For each \(m\)-ary predicate \(P\), define
\[
 P^{\mathcal M}
   =\{(a_1,\ldots,a_m)\in V^m:
                    P(a_1,\ldots,a_m)\in\Gamma_\infty\}.
\]
When \(m=0\), this gives the corresponding truth value by membership
of the nullary atom.  The domain is nonempty.
We prove simultaneously, by formula complexity, that every formula
in \(\Gamma_\infty\) is true at \((\mathcal M,s_0)\) and every
formula in \(\Delta_\infty\) is false there.

The atomic claims follow from the definition of the relations and
the absence of an overlap.  Bottom is always false and is absent
from the antecedent.  For an antecedent conjunction, both conjuncts
are present in \(\Gamma_\infty\); for a succedent conjunction, one
conjunct is present in \(\Delta_\infty\).  The two disjunction cases
are dual.  An antecedent implication has either its antecedent in
\(\Delta_\infty\) or its consequent in \(\Gamma_\infty\), making it
true.  A succedent implication has its antecedent in \(\Gamma_\infty\)
and consequent in \(\Delta_\infty\), making it false.

For \(\forall xA\in\Gamma_\infty\), every \(A[a/x]\), \(a\in V\),
belongs to \(\Gamma_\infty\).  The induction hypothesis and the
substitution lemma give truth of \(A\) under every value for \(x\),
since every domain element is one of these variables.
For \(\exists xA\in\Delta_\infty\), the same argument with every
instance in \(\Delta_\infty\) makes the existential false.
For \(\exists xA\in\Gamma_\infty\), its recorded witness instance
is true by induction, making the existential true.
For \(\forall xA\in\Delta_\infty\), its recorded witness instance
is false, making the universal false.
Substitution of a variable changes no logical complexity, so every
instance used here lies within the induction hypothesis.

The resulting structure and assignment falsify the original sequent.
Thus a valid finite sequent cannot be unprovable.  Together with
soundness, this proves the stated equivalence.
\end{proof}
\end{annexitem}

\section{Craig interpolation in general signatures}

\begin{annexitem}{annex:p7.interp.3}{thm:craig-interpolation}{Craig interpolation}
\begin{proof}
\paragraph{The relational case.}
Let \(A,B\) be equality-free purely relational sentences with
\(\models A\to B\).
Theorem~\ref{thm:p07-relational-cut-free-completeness} supplies a
cut-free derivation of \(A\Rightarrow B\).
Put \(A\) in the first antecedent part and \(B\) in the second
succedent part of its partition.
Theorem~\ref{thm:p07-maehara-partition} supplies a sentence \(J\)
in the common predicate vocabulary, with derivations of
\(A\Rightarrow J\) and \(J\Rightarrow B\).
Soundness gives \(\models A\to J\) and \(\models J\to B\).
This establishes the relational case, including an empty common
predicate vocabulary.

\paragraph{A finite relational signature for each side.}
Now let \(A,B\) be arbitrary first-order sentences with
\(\models A\to B\).
Write \(\Sigma_1,\Sigma_2\) for their respective finite sets of
nonlogical symbols, with arities retained.
Only the union of these finite signatures is needed.
Every nonempty structure in that union extends to any unused
original symbols by assigning empty predicates and constant-valued
functions at one fixed element.

Introduce one fresh binary predicate \(E\).
For every \(m\)-ary function symbol \(f\), introduce a distinct
\((m+1)\)-ary predicate \(G_f\).  Thus a constant \(c\), regarded
as a function of arity zero, has a unary predicate \(G_c\).
Keep original predicate symbols unchanged.
The translated language is purely relational and has no logical
equality, functions, or constants.

For a finite signature \(\Sigma\), let \(\operatorname{Ax}_\Sigma\)
be the conjunction of the following finite set of relational sentences.
First require \(E\) to be an equivalence relation:
\[
\begin{gathered}
 \forall x\,E(x,x),\\
 \forall x\,y\,(E(x,y)\to E(y,x)),\\
 \forall x\,y\,z\,
       ((E(x,y)\land E(y,z))\to E(x,z)).
\end{gathered}
\]
For every original predicate in \(\Sigma\) and every graph predicate
associated with a function in \(\Sigma\), require congruence in all
coordinates.  For such a predicate \(H\) of arity \(r\geq1\), this is
\[
 \forall\bar x\,\bar y\,
 \left[
  \left(\bigwedge_{i=1}^{r}E(x_i,y_i)\right)
  \to\bigl(H(\bar x)\leftrightarrow H(\bar y)\bigr)
 \right].
\]
For arity zero, congruence is a tautology and is omitted.
Finally, for each \(m\)-ary function \(f\in\Sigma\), include
\[
\begin{aligned}
 &\forall x_1\cdots x_m\,\exists z\,
                   G_f(x_1,\ldots,x_m,z),\\
 &\forall x_1\cdots x_m\,z\,w\,
 \bigl((G_f(x_1,\ldots,x_m,z)\land
        G_f(x_1,\ldots,x_m,w))\to E(z,w)\bigr).
\end{aligned}
\]
For \(m=0\), the first sentence is \(\exists z\,G_c(z)\), and the
second is uniqueness modulo \(E\) for that unary graph.
The equivalence and conjunction symbols in these axioms have their
ordinary definitions in the selected primitives.
Only \(E\), original predicates from \(\Sigma\), and graphs of
functions from \(\Sigma\) occur in \(\operatorname{Ax}_\Sigma\).

\paragraph{Terms and formulas.}
For a term \(t\), define a relational formula \(V_t(z)\), where the
output variable \(z\) is fresh for \(t\), by
\[
\begin{aligned}
 V_x(z)&=E(x,z),\\
 V_{f(t_1,\ldots,t_m)}(z)
   &=\exists z_1\cdots z_m
       \left(\bigwedge_{i=1}^{m}V_{t_i}(z_i)
                       \land G_f(z_1,\ldots,z_m,z)\right).
\end{aligned}
\]
All auxiliary variables are distinct and fresh for the entire input
formula, and bound variables are renamed when necessary.
For a constant this gives \(V_c(z)=G_c(z)\).
The empty conjunction is \(\top\), and an empty existential block
does nothing.

The formula translation \((\cdot)^r\) fixes \(\bot\), commutes with
the Boolean connectives and quantifiers, and translates atoms by
\[
\begin{aligned}
 (P(t_1,\ldots,t_m))^r
  &=\exists z_1\cdots z_m
       \left(\bigwedge_{i=1}^{m}V_{t_i}(z_i)
                       \land P(z_1,\ldots,z_m)\right),\\
 (t=u)^r
  &=\exists z\,w\,
       (V_t(z)\land V_u(w)\land E(z,w)).
\end{aligned}
\]
A nullary atom is unchanged.
Every introduced variable is bound.  No original symbol is replaced
by the graph of a different symbol, so the translation preserves
which original side owns each predicate, function, or constant.

\paragraph{The quotient truth lemma.}
Let \(N\) be a nonempty relational structure satisfying
\(\operatorname{Ax}_\Sigma\), with domain \(D\).
Write \([a]\) for the \(E^N\)-class of \(a\in D\), and let \(Q\)
have domain \(D/E^N\).
Interpret an original predicate by
\[
 P^Q([a_1],\ldots,[a_m])
       \quad\Longleftrightarrow\quad N\models P(a_1,\ldots,a_m).
\]
Its congruence axiom makes this independent of the representatives.
Interpret \(f^Q([a_1],\ldots,[a_m])\) as the unique class \([b]\)
such that \(N\models G_f(a_1,\ldots,a_m,b)\).
Totality supplies a class, uniqueness makes it unique, and input
congruence makes it independent of the representatives of the inputs.
For constants the same clauses give a unique output class.
Logical equality in \(Q\) is identity of equivalence classes.

If \(s\) is an assignment into \(D\), write \(\bar s(x)=[s(x)]\).
Induction on a term proves
\[
 N,s[z:=a]\models V_t(z)
 \quad\Longleftrightarrow\quad
 [a]=t^Q(\bar s).
\]
For a variable this is exactly the interpretation of \(E\).
For a function term, witnesses to the translated subterms give their
correct argument classes by induction, so its graph atom gives the
stated value class.  Conversely, choose representatives of the finitely
many argument classes.  Their term formulas hold by induction.
Totality gives a graph output in the required value class, and
congruence in the output coordinate permits replacement by the
specified representative \(a\).  This supplies all existential witnesses.
The constant case uses the same totality, uniqueness, and output
congruence with no arguments.

Induction on formulas now gives
\[
 N,s\models C^r
 \quad\Longleftrightarrow\quad Q,\bar s\models C.
\]
For an original predicate atom, the term lemma and predicate congruence
give both directions.  For equality, the term lemma identifies the
two value classes, and \(E\) says exactly that they agree.
Bottom and the Boolean connectives follow from their truth clauses.
For a quantifier, every value in \(D\) maps to an equivalence class,
and every class has a representative in \(D\).  Thus existential
witnesses lift and descend, and the universal quantifier ranges
over all classes exactly as required.  This proves the full truth lemma.

\paragraph{The interpolant uses only common symbols.}
Set \(\operatorname{Ax}_i=\operatorname{Ax}_{\Sigma_i}\).
We have the relational validity
\[
 \models(\operatorname{Ax}_1\land A^r)
                 \to(\operatorname{Ax}_2\to B^r).
\]
Indeed, a countermodel would satisfy both axiom conjunctions,
\(A^r\), and \(\neg B^r\).
Its quotient for \(\Sigma_1\cup\Sigma_2\) would satisfy
\(A\land\neg B\) by the truth lemma, contradicting the original
valid implication.
Apply the already proved relational case to these two sentences.
It gives a relational sentence \(J\) with
\[
\begin{aligned}
 &\models(\operatorname{Ax}_1\land A^r)\to J,\\
 &\models J\to(\operatorname{Ax}_2\to B^r).
\end{aligned}
\]
Its predicate vocabulary is contained in the intersection of the
two translated vocabularies.  Hence it can contain \(E\), common
original predicates, and \(G_f\) only for a function or constant
\(f\) occurring in both original sentences.

Decode \(J\) into a first-order sentence \(\theta\) by the atomic
replacements
\[
 E(x,y)\mapsto x=y,\qquad
 G_f(x_1,\ldots,x_m,z)\mapsto f(x_1,\ldots,x_m)=z,
\]
keeping original predicates and commuting with all logical operations.
For a constant the second replacement is \(G_c(z)\mapsto c=z\).
This introduces only common original symbols; equality is logical,
as permitted in the theorem.

For any original structure \(M\), its canonical relational expansion
\(M^r\) has the same domain, interprets \(E\) as actual equality,
and interprets each \(G_f\) as the graph of \(f^M\).
It satisfies both axiom conjunctions.
Its quotient is canonically isomorphic to \(M\), so the quotient
truth lemma yields \(M\models C\) iff \(M^r\models C^r\).
Also induction directly on \(J\) yields
\(M^r\models J\) iff \(M\models\theta\).
The first interpolant implication therefore gives
\(\models A\to\theta\); the second, since \(M^r\models
\operatorname{Ax}_2\), gives \(\models\theta\to B\).
This proves sentence interpolation in the general signature.

\paragraph{Free variables.}
Let now \(\varphi,\psi\) be formulas.
For each variable \(x\in\FV(\varphi)\cup\FV(\psi)\), choose a distinct
fresh constant symbol \(c_x\), and replace its free occurrences by
\(c_x\).  Let the resulting sentences be \(\varphi^c,\psi^c\).
Validity of \(\varphi\to\psi\) under every assignment implies
validity of \(\varphi^c\to\psi^c\) under every interpretation of the
new constants.
Sentence interpolation produces \(\theta^c\) with only common
symbols.  A constant \(c_x\) can therefore occur in it only when
\(x\in\FV(\varphi)\cap\FV(\psi)\).

Rename bound variables of \(\theta^c\) away from this finite set,
then replace each \(c_x\) by \(x\), including occurrences inside terms.
Call the result \(\theta\).
For any original structure and assignment \(s\), interpret
\(c_x\) as \(s(x)\).  Induction on terms and formulas identifies
the truth of \(\theta^c\) in this expansion with the truth of
\(\theta\) at \(s\), and likewise for the two input formulas.
The two sentence implications thus become
\(\models\varphi\to\theta\) and \(\models\theta\to\psi\).
The construction gives the asserted free-variable intersection.
Distinct new constant symbols were never required to denote
distinct elements, so this argument covers assignments with equal
values as well.
\end{proof}
\end{annexitem}

\section{Beth definability}

\begin{annexitem}{annex:p7.interp.4}{thm:beth}{Beth definability}
\begin{proof}
Write \(L\) for the signature obtained by removing \(R\), and let
\(n\geq0\) be the arity of \(R\).
Thus \(T\subseteq\operatorname{Sent}(L\cup\{R\})\).
Choose a fresh \(n\)-ary predicate \(R'\), and let \(T'\) be obtained
by uniformly replacing \(R\) by \(R'\) in every sentence of \(T\).
All symbols of \(L\), including functions and constants, remain fixed.
Choose fresh distinct constant symbols \(c_1,\ldots,c_n\).
If \(n=0\), this list is empty.

\paragraph{One finite interpolation problem.}
The set of sentences
\[
 T\cup T'\cup\{R(\bar c),\neg R'(\bar c)\}
\]
is unsatisfiable.  Otherwise its model would give two expansions of
one \(L\)-reduct satisfying \(T\): the first uses its interpretation
of \(R\), and the second uses its interpretation of \(R'\) as \(R\).
Implicit definability makes these relations equal, contradicting the
two displayed literals.
Compactness, Theorem~\ref{thm:fol-compactness}, therefore supplies
finite \(\Delta,\Gamma\subseteq T\) for which
\[
 \Delta\cup\Gamma'\cup
       \{R(\bar c),\neg R'(\bar c)\}
\]
is unsatisfiable.  If a literal was absent from the initial finite
unsatisfiable subset, adding it preserves unsatisfiability.
The primes on \(\Gamma'\) denote the same uniform predicate renaming.
Set
\[
 A=\left(\bigwedge\Delta\right)\land R(\bar c),
 \qquad
 B=\left(\bigwedge\Gamma'\right)\to R'(\bar c),
\]
with empty conjunctions equal to \(\top\).
The finite unsatisfiability statement is exactly \(\models A\to B\).

Craig interpolation, Theorem~\ref{thm:craig-interpolation}, gives a
sentence \(J(\bar c)\) with
\[
 \models A\to J(\bar c),\qquad
 \models J(\bar c)\to B.
\]
Here the notation \(J(\bar c)\) names its possible parameter constants,
not free variables.
The predicate \(R\) occurs only on the first side, and \(R'\) only
on the second.  Neither can occur in \(J\).
Apart from the fresh constants, every nonlogical symbol in \(J\)
lies in \(L\), indeed in the intersection of the finite \(L\)-vocabularies
of the two conjunctions.
The two implications can be written
\[
\begin{aligned}
 &\Delta\models R(\bar c)\to J(\bar c),\\
 &\Gamma'\models J(\bar c)\to R'(\bar c).
\end{aligned}
\]
Uniformly rename \(R'\) back to \(R\) in the second implication.
Since \(J\) contains neither predicate, it remains unchanged.
Thus, for the single finite set \(T_0=\Delta\cup\Gamma\subseteq T\),
\[
 T_0\models R(\bar c)\leftrightarrow J(\bar c).
\]

\paragraph{Removing the parameters.}
Take pairwise distinct variables \(x_1,\ldots,x_n\).
First rename the bound variables of \(J\) away from them, then
replace each fresh constant \(c_i\) by \(x_i\), including occurrences
inside compound terms.  Call the result \(\theta(\bar x)\).
It is an \(L\)-formula with no free variables outside \(\bar x\).
Its nonlogical vocabulary is finite, as every formula is finite.

Let \(M\models T_0\), and choose any tuple
\((a_1,\ldots,a_n)\) from its domain.
Expand \(M\) by interpreting \(c_i\) as \(a_i\).
The constants do not occur in \(T_0\), so this expansion still
satisfies \(T_0\).  The last displayed consequence gives
\[
 R^M(a_1,\ldots,a_n)
 \quad\Longleftrightarrow\quad
 M,\bar x\mapsto\bar a\models\theta(\bar x),
\]
where the identification on the right follows by the ordinary
term and formula substitution induction.
Distinct constant symbols may denote the same element; no inequality
between tuple entries has been assumed.
Since the tuple was arbitrary,
\[
 T_0\models\forall\bar x\,
               (R(\bar x)\leftrightarrow\theta(\bar x)).
\]
First-order completeness, Theorem~\ref{thm:fol-completeness}, yields
the displayed derivation from \(T_0\), and therefore from \(T\).
The same finite support works for all models and all tuples:
\(\Delta,\Gamma\), and \(J\) were chosen before any model or assignment.

When \(n=0\), the argument uses no fresh constants or variables and
gives the sentence \(R\leftrightarrow\theta\) directly.
If \(T\) is inconsistent, compactness gives a finite inconsistent
subset and \(\theta=\bot\) also serves as an explicit definition.
The theorem makes no assumption that every \(L\)-structure has a
\(T\)-expansion; it compares those expansions that satisfy \(T\).
\end{proof}
\end{annexitem}
